\subsection*{Questão 05}

\begin{questionbox}{05}
  Para a formação de uma equipe de desenvolvimento de um novo sistema de Inteligência Artificial (IA), um gerente de projetos precisa selecionar 4 especialistas de um grupo de 9. 

  No entanto, o grupo contém 3 especialistas que, por terem visões de arquitetura conflitantes, não podem trabalhar juntos no mesmo projeto. 

  De quantas maneiras o gerente pode formar uma equipe de 4 especialistas sem formar qualquer par entre os 3 citados?
\end{questionbox}

\begin{solutionbox}
  Sejam os 3 especialistas conflitantes. A restrição implica que não podemos escolher dois ou mais deles simultaneamente. Logo, podemos escolher no máximo 1 desses 3.

  Dividimos em casos:

  \textbf{Caso 1:} Nenhum dos 3 é escolhido.

  Selecionamos os 4 membros apenas dentre os outros 6:
  \[
  \binom{6}{4} = 15.
  \]

  \textbf{Caso 2:} Exatamente 1 dos 3 é escolhido.

  Escolhemos 1 dos 3 e os outros 3 dentre os 6 restantes:
  \[
  \binom{3}{1} \cdot \binom{6}{3} = 3 \cdot 20 = 60.
  \]
\end{solutionbox}

\newpage

\begin{solutionbox}
  Somando os casos:
  \[
  15 + 60 = 75.
  \]

  Portanto, o número de equipes possíveis é:
  \[
  \boxed{75}.
  \]
\end{solutionbox}